\section{Introduction}
\label{sec:introduction}

The Heavy Photon Search (HPS) experiment performs a dark-sector-motivated search for a
narrow, electroproduced prompt $A^\prime \rightarrow e^+e^-$ resonance in the
reconstructed invariant-mass spectrum. In this class of analysis the central
statistical problem is the control of the smooth QED background immediately surrounding
each test mass. The analysis documented here addresses that problem with a global
Gaussian-process-regression (GPR) model for the invariant-mass spectrum of each active
dataset, a profiled multi-bin likelihood for local signal extraction, \CLs\ upper
limits in signal-yield and $\eps^2$ space, and a statistically consistent combination
of multiple datasets in a shared coupling parameter
\cite{HPS2015DarkPhoton,HPS2016PromptLong,HPSExperiment2022}.

This note is written from the current analysis implementation and from the mature
physics logic of the study rather than from the earliest exploratory drafts. In
particular, the baseline local-significance calculation is the profiled Gaussian
multi-bin likelihood-ratio test rather than the older blind-window sum-count
approximation, and the signal template is treated as a full Gaussian line shape whose
blinded-bin slice enters the likelihood without renormalization. That convention keeps
the fitted amplitude aligned with the physically meaningful quantity, namely the total
local signal yield of a narrow resonance.

\subsection{Published HPS prompt-resonance benchmarks}

\subsubsection{2015 engineering run}

The published 2015 HPS prompt search analyzed \SI{1170}{nb^{-1}} collected at a beam
energy of \SI{1.056}{GeV} and searched for a resonance in the invariant-mass interval
\SIrange{19}{81}{MeV} \cite{HPS2015DarkPhoton}. The resonance fit was performed in a
sliding local window stepped in \SI{0.5}{MeV} increments. The fit window was taken to
be $14\,\sigma_m$ wide below \SI{39}{MeV} and $13\,\sigma_m$ wide above
\SI{39}{MeV}, with the background described by a fifth-order Chebyshev polynomial below
\SI{39}{MeV} and a third-order Chebyshev polynomial above that transition. The most
significant excess occurred at \SI{37.7}{MeV} with a local $p$-value of
$1.7\times10^{-3}$; after the look-elsewhere correction the global $p$-value was 17\%.
Since no significant signal was found, the analysis set 95\% power-constrained limits
excluding couplings down to the level of $\eps^2 \simeq 6\times10^{-6}$.

\subsubsection{2016 engineering run}

The published 2016 prompt search analyzed the full 2016 engineering-run sample,
\SI{10608}{nb^{-1}} at \SI{2.3}{GeV}, and searched the invariant-mass interval
\SIrange{39}{179}{MeV} \cite{HPS2016PromptLong}. In the prompt channel the smallest
local $p$-value was $6.38\times10^{-3}$ at \SI{94}{MeV}; after accounting for the
look-elsewhere effect the global significance was approximately $1\,\sigma$. The
corresponding 95\% \CLs\ upper limits excluded couplings of order $10^{-5}$, with
68\% and 95\% expected bands reported explicitly and systematic uncertainties from the
mass resolution and radiative-fraction determination propagated into the final result.

The 2016 publication also sharpened the physics interpretation of the prompt search in
ways that are directly relevant for the present note. The conversion from signal yield
to $\eps^2$ was written in terms of the radiative trident rate, the radiative fraction
was determined from selected Monte Carlo samples of radiative tridents, tridents, and
converted wide-angle Bremsstrahlung, and the mass-resolution systematics were carried
through to the prompt-significance and limit curves. Those ingredients remain central
to the present GPR-based analysis even though the background model itself is now
different.

\begin{figure}[H]
\centering
\begin{subfigure}[t]{0.32\linewidth}
  \centering
  \graphicorplaceholder{\linewidth}{published_reference_figs/hps2015_published_limit.png}
  \caption{Published 2015 95\% power-constrained limit, reaching
  $\eps^2 \sim 6\times10^{-6}$ over \SIrange{19}{81}{MeV}.}
\end{subfigure}
\hfill
\begin{subfigure}[t]{0.32\linewidth}
  \centering
  \graphicorplaceholder{\linewidth}{published_reference_figs/hps2016_published_prompt_pvalue.png}
  \caption{Published 2016 prompt-search local/global $p$-value scan over
  \SIrange{39}{179}{MeV}.}
\end{subfigure}
\hfill
\begin{subfigure}[t]{0.32\linewidth}
  \centering
  \graphicorplaceholder{\linewidth}{published_reference_figs/hps2016_published_prompt_limit.png}
  \caption{Published 2016 prompt-search 95\% \CLs\ limit, excluding couplings of order
  $10^{-5}$.}
\end{subfigure}
\caption{Published HPS prompt-resonance benchmarks. The 2015 result established the
first HPS prompt exclusion in the \SIrange{19}{81}{MeV} region, while the 2016 result
extended the reach to higher masses and incorporated a more mature treatment of
systematic uncertainties and expected-limit bands.}
\label{fig:published-hps-context}
\end{figure}

\subsection{Motivation for the present GPR analysis}

The motivation for the current analysis is not simply to replace one smooth-background
fit with another. The earlier local-window resonance searches were necessarily tied to
window-dependent background parameterizations and per-window stability choices. In the
current study, the global GP description of the invariant-mass spectrum provides a
single mean-and-covariance prediction for the blinded bins at each mass point. That
allows the signal extraction, local significance, expected limits, and shared-$\eps^2$
combination to be formulated within one coherent statistical framework.

From an analysis perspective, the most important practical gain is stability. The GP
approach is demonstrably smoother than the earlier piecewise background fits and is
therefore better suited to extended search windows and to the combination of data sets
with different resolutions, exposures, and kinematic coverage. That stability is what
makes it realistic to validate wider search intervals, to compare 2015 and 2016 within
a common framework, and eventually to incorporate the 2021 data sample in the same
statistical language.

\subsection{Scope of this note}

The note has two aims. First, it documents the current narrow prompt
$e^+e^-$-resonance analysis at a level suitable for collaboration review and later
adaptation into a journal manuscript. Second, it organizes the available figures,
tables, and placeholders according to the final analysis sequence so that the completed
results can be inserted without restructuring the document. Where a 2021 input or a
post-rerun study product is not yet finalized, the corresponding text is kept explicit
about that status rather than written as if the analysis were already complete.
